\section{Three-term swap codes and one- and two-sided interval-mask bounds}\label{sec:swapcodes}

\subsection{Three-term swap codes exceeding \texorpdfstring{$d(d+1)/2$}{d(d+1)/2} at \texorpdfstring{$d \in \{6, 7, 9, 12\}$}{d in \{6, 7, 9, 12\}} and over all fields at \texorpdfstring{$d = 4$}{d = 4}}\label{sec:remark-5-6}

For finite matrix families $A, B : \alpha \to \Mat_{d\times d}(F)$ over a field $F$, the swap code
(\leanname{SwapRigidity.swapCode}) is the solution space
\[
  \leanname{SwapRigidity.swapCode}(A, B)
  \;:=\;
  \Bigl\{X \in \Mat_{d\times d}(F) \;:\; \sum_{a \in \alpha} A_a X B_a^{\T} = X^{\T}\Bigr\}.
\]
In \cite[Theorem~5.1 and Corollary~5.5]{SeventeenthPaper2026} (where the coordinate form of
Theorem~5.1 is Lean-verified and Corollary~5.5 is proved in text), any two-term
swap code ($\abs{\alpha} = 2$) has dimension at most the symmetric pairing value $d(d+1)/2$, proved
by triangularizing a matrix pencil over an algebraic closure. Immediately afterward, Remark~5.6 of
\cite{SeventeenthPaper2026} noted that three terms can exceed $d(d+1)/2$, giving an
explicit three-term family of dimension $8 > 6$ at $d = 3$ over every field and an example of
dimension $11 > 10$ at $d = 4$ (checked numerically there over $\F_2, \F_3, \F_5, \F_7$),
and remarked: \emph{``We have no three-term example for $d \ge 5$.''}

Here we give explicit three-term examples above $d(d+1)/2$ at $d \in \{6, 7, 9, 12\}$ over every
field $F$, together with a proof over every field at $d = 4$ (for a padded family matching the
dimension $11 > 10$ of \cite[Remark~5.6]{SeventeenthPaper2026}, previously checked over
$\F_2, \F_3, \F_5, \F_7$), by lifting the $3 \times 3$ skew-pair construction into $m \times m$
blocks (\lean{threeTermBlockA}, \lean{threeTermBlockB}) and adjoining a diagonal tail
(\lean{threeTermPaddedA}, \lean{threeTermPaddedB}). Write the three off-diagonal pairs of $\{0,1,2\}$
as $(i_0,j_0) = (0,1)$, $(i_1,j_1) = (0,2)$, and $(i_2,j_2) = (1,2)$, and set
\[
  A_t \;:=\; (E_{i_t, j_t} - E_{j_t, i_t}) \otimes I_m,
  \qquad
  B_t \;:=\; -A_t \;=\; (E_{j_t, i_t} - E_{i_t, j_t}) \otimes I_m
  \qquad (t \in \{0,1,2\}).
\]
Because $E_{j_t, i_t}^{\T} - E_{i_t, j_t}^{\T} = E_{i_t, j_t} - E_{j_t, i_t}$, we have
$B_t^{\T} = A_t$, so the swap-code equation $\sum_{t=0}^2 A_t X B_t^{\T} = X^{\T}$ is
\[
  \sum_{t=0}^2 \bigl((E_{i_t, j_t} - E_{j_t, i_t}) \otimes I_m\bigr)\,X\,
    \bigl((E_{i_t, j_t} - E_{j_t, i_t}) \otimes I_m\bigr)
  \;=\;
  X^{\T}.
\]

\begin{theorem}[Dimension $4m(m+1)$ of the three-term block family on $d = 3m$]\label{thm:three-term-kronecker}
\formal{SwapCodeAnswersProofs.lean}{bonus1_ThreeTermBlockSwapCodeDim}
Let $F$ be any field and let $m \ge 1$. Partition $X \in \Mat_{3m\times 3m}(F)$ into a $3 \times 3$
grid of $m \times m$ blocks $(X_{i,j})_{0 \le i,j < 3}$.
\begin{enumerate}[label=(\alph*)]
\item $X$ lies in $\leanname{SwapRigidity.swapCode}(A, B) = \{X : \sum_{t=0}^2 A_t X B_t^{\T} = X^{\T}\}$
  if and only if every block $X_{i,j}$ is symmetric ($X_{i,j}^{\T} = X_{i,j}$ for all
  $0 \le i, j < 3$) and the three diagonal blocks sum to zero ($X_{0,0} + X_{1,1} + X_{2,2} = 0$;
  \lean{IsThreeByThreeSymmetricZeroTraceBlocks}).
\item Consequently, over every field $F$,
  \[
    \dim_F \leanname{SwapRigidity.swapCode}(A, B) \;=\; 8 \cdot \frac{m(m+1)}{2} \;=\; 4m(m+1).
  \]
  Since $4m(m+1) - \frac{3m(3m+1)}{2} = \frac{m(5 - m)}{2}$, this exceeds the pairing value
  $d(d+1)/2$ ($d = 3m$) for $m \in \{1, 2, 3, 4\}$ ($d = 3: 8 > 6$, matching
  \cite[Remark~5.6]{SeventeenthPaper2026}; $d = 6: 24 > 21$; $d = 9: 48 > 45$; and
  $d = 12: 80 > 78$).
\end{enumerate}
\end{theorem}
\begin{outline}
For each $t \in \{0,1,2\}$, multiplying on both sides by $(E_{i_t,j_t} - E_{j_t,i_t}) \otimes I_m$
places $+X_{j_t, i_t}$ in block $(i_t, j_t)$, $+X_{i_t, j_t}$ in block $(j_t, i_t)$, $-X_{j_t, j_t}$
in block $(i_t, i_t)$, and $-X_{i_t, i_t}$ in block $(j_t, j_t)$. Summing over the three pairs sends
each off-diagonal block $(i, j)$ ($i \ne j$) to $+X_{j,i}$ and each diagonal block $(i, i)$ to
$-\sum_{j \ne i} X_{j,j}$. Equating with $(X_{j,i}^{\T})_{i,j}$ gives $X_{j,i} = X_{j,i}^{\T}$ for
$i \ne j$ together with $-X_{1,1} - X_{2,2} = X_{0,0}^{\T}$, $-X_{0,0} - X_{2,2} = X_{1,1}^{\T}$,
and $-X_{0,0} - X_{1,1} = X_{2,2}^{\T}$. Over every ring (including characteristic~$2$), subtracting
pairs of diagonal equations yields $X_{0,0}^{\T} - X_{1,1}^{\T} = X_{0,0} - X_{1,1}$ and cyclic
analogues, which together with $X_{2,2} = -X_{0,0}^{\T} - X_{1,1}$ is equivalent to
$X_{i,i}^{\T} = X_{i,i}$ for $i \in \{0,1,2\}$ and $X_{0,0} + X_{1,1} + X_{2,2} = 0$. Eight
independent symmetric $m \times m$ blocks ($6$ off-diagonal and $2$ free diagonal blocks) each have
dimension $m(m+1)/2$, giving $4m(m+1)$.
\end{outline}

Embedding the $3m \times 3m$ family in the top-left corner of $\Mat_{d\times d}(F)$ ($d = 3m + s$)
and setting both $A_2$ and $B_2$ equal to $I_s$ on the trailing $s \times s$ diagonal block extends
the construction to $d = 4$ and $d = 7$.

\begin{theorem}[Padded three-term swap codes across $d \in \{3, 4, 6, 7, 9, 12\}$]\label{thm:three-term-unified}
\formal{SwapCodeAnswersProofs.lean}{bonus1_ThreeTermPaddedSevenAndFive}
\formal{DualityAndSwapConsequences.lean}{bonus1_UnifiedThreeTermExceedsPairingBound}
Over every field $F$:
\begin{enumerate}[label=(\alph*)]
\item Padding $m = 2$ ($3m = 6$) to $d = 7$ achieves dimension $29 > 28 = \frac{7 \cdot 8}{2}$;
  padding $m = 1$ ($3m = 3$) to $d = 4$ achieves $11 > 10 = \frac{4 \cdot 5}{2}$; and padding
  $m = 1$ to $d = 5$ achieves $15 = \frac{5 \cdot 6}{2}$.
\item Consequently, for every field $F$ and every $d \in \{3, 4, 6, 7, 9, 12\}$, there exist
  $A, B : \mathrm{Fin}\,3 \to \Mat_{d\times d}(F)$ with
  \[
    \dim_F \leanname{SwapRigidity.swapCode}(A, B) \;>\; \frac{d(d+1)}{2}.
  \]
\end{enumerate}
\end{theorem}
\begin{outline}
On the trailing $s = d - 3m$ coordinates, $A_2$ and $B_2^{\T}$ both act as $+I_s$, while $A_0, B_0,
A_1, B_1$ vanish outside the top-left $3m \times 3m$ block. Inspecting the coupled cross-blocks
between the active $(1,2)$ pair and the $s$-coordinate tail yields $24 + 5 = 29$ dimensions at
$(m, d) = (2, 7)$, $8 + 3 = 11$ at $(1, 4)$, and $8 + 7 = 15$ at $(1, 5)$. Combining $(1,4)$ and
$(2,7)$ with $m \in \{1, 2, 3, 4\}$ in \Cref{thm:three-term-kronecker} covers
$d \in \{3, 4, 6, 7, 9, 12\}$.
\end{outline}

More generally (proved in text; not formalized for general $m, s$), for $d = 3m + s$, partitioning
$X \in \Mat_{d\times d}(F)$ into three $m$-coordinate blocks $0, 1, 2$ and a trailing
$s$-coordinate block $3$, the top-left $3m \times 3m$ blocks still contribute $4m(m+1)$ dimensions by
\Cref{thm:three-term-kronecker}, while equating $\sum_{t=0}^2 A_t X B_t^{\T} = X^{\T}$ on row/column~$3$
gives $X_{0,3} = 0$, $X_{3,0} = 0$, $X_{2,3} = X_{3,1}^{\T}$ ($ms$ free entries),
$-X_{1,3} = X_{3,2}^{\T}$ ($ms$ free entries), and $X_{3,3} = X_{3,3}^{\T}$ ($s(s+1)/2$ free entries).
Thus for $m = \lfloor d/3\rfloor$ and $s = d - 3m \in \{0, 1, 2\}$, this block-and-tail construction
achieves dimension
$4m(m+1) + \frac{s(s+1)}{2} + 2ms = \frac{d(d+1)}{2} + \frac{m(5 - m - 2s)}{2}$, which is at most
$d(d+1)/2$ for every $(m \ge 1, s \ge 0)$ with $d \ge 5$ outside $\{6, 7, 9, 12\}$ (giving
$15 = 15$ at $d = 5$, $35 < 36$ at $d = 8$, $55 = 55$ at $d = 10$, $63 < 66$ at $d = 11$,
$120 = 120$ at $d = 15$, and strictly less than $d(d+1)/2$ once $m \ge 6$). Whether a three-term swap
code can exceed $d(d+1)/2$ remains open for $d = 5$ and for every $d \ge 8$ outside $\{9, 12\}$.

\subsection{One-sided multi-scale and two-sided two-scale (coarse--fine) bounds (Problem~9.3)}\label{sec:problem-9-3}

Problem~9.3 of \cite{SeventeenthPaper2026} asks for upper bounds on
$\dim_F \leanname{SwapRigidity.swapCode}(A, B)$ when $d = 2^m$ and both $A_a$ and $B_a$ lie in the
span of dyadic interval masks across all scales composed with cyclic rotations
(\cite[Corollary~6.13]{SeventeenthPaper2026}, proved in text, covers one-sided tensor-factor
rules commuting with a single non-derogatory matrix on one address digit). We first record the
one-sided endpoint-refinement baseline across all scales, then prove a characteristic-free two-sided
bound for two-scale (coarse--fine) interval families that improves the leading coefficient over the
one-sided Theorem~TP baseline $\bigl(d^2 + \min(s_A + c_A, s_B + c_B)d\bigr)/2$ implied by the
decomposition counts alone (where each element of $S_A, S_B$ may carry its own singleton mask, as in
\Cref{cor:two-sided-win}) when $s_A/d$ and $s_B/d$ are bounded away from $0$ and $1$ in $(0, 1)$
and the run and coarse counts are
$o(d)$, and finally contrast interval masks with periodic Walsh digit masks.

For contiguous cyclic or aligned dyadic intervals (\lean{IsCyclicInterval},
\lean{IsAlignedDyadicInterval}; spans \lean{InMultiscaleIntervalRotSpan},
\lean{InAlignedDyadicIntervalRotSpan}), each interval mask introduces at most two boundary
transitions around $\Z_d$, so applying Theorem~TP (\cite[Theorem~6.4]{SeventeenthPaper2026},
Lean-verified) to the
common refinement of the interval endpoints on one side at a time gives a nontrivial upper bound
whenever at least one side uses $o(d)$ intervals in total:

\begin{proposition}[One-sided endpoint-refinement bound for multi-scale interval masks]\label{prop:multiscale-interval}
\formal{SwapCodeAnswersProofs.lean}{bonus2_MultiscaleIntervalBothSides,bonus2_MultiscaleFullProfileAttainsFullSwap}
\formal{DualityAndSwapConsequences.lean}{bonus2_SublinearProfileAndTendstoHalf,bonus2_SublinearProfileExplicitThreshold}
Let $F$ be any field, $m \ge 0$, and $d = 2^m$.
\begin{enumerate}[label=(\alph*)]
\item Every \lean{InAlignedDyadicIntervalRotSpan} family lies in \lean{InMultiscaleIntervalRotSpan},
  and whenever finite families $A, B : \alpha \to \Mat_{d\times d}(F)$ lie in
  \lean{InMultiscaleIntervalRotSpan} with per-scale interval counts $M_A, M_B : \mathrm{Fin}\,m \to \N$,
  setting $\beta_A := \max\bigl(1,\, 2\sum_j M_A(j)\bigr)$ and
  $\beta_B := \max\bigl(1,\, 2\sum_j M_B(j)\bigr)$,
  \[
    2\dim_F \leanname{SwapRigidity.swapCode}(A, B)
    \;\le\;
    d^2 + d\min(\beta_A, \beta_B).
  \]
\item Along any sequence of fields $(F_m)$ and families $(A_m, B_m)$ whose per-scale interval sums
  $S_{A,m} := \sum_j M_{A,m}(j)$ and $S_{B,m} := \sum_j M_{B,m}(j)$ satisfy
  $\min(S_{A,m}, S_{B,m}) / 2^m \to 0$ as $m \to \infty$, for every $\eps > 0$ one eventually has
  $\dim_{F_m}{}$\leanname{SwapRigidity.swapCode}$(A_m, B_m) \le (\frac{1}{2} + \eps)d^2$ (uniformly
  for all $m \ge \lfloor 1/\eps\rfloor + 1$ whenever $\min(S_{A,m}, S_{B,m}) \le (\eps/2)\,2^m$).
\item At the full aligned dyadic profile $M_{\mathrm{full}}(j) = 2^{j+1}$ (satisfying
  $\sum_{j=0}^{m-1} 2^{j+1} \ge d$ for $m \ge 1$), there exist aligned dyadic interval-rotation
  families attaining $\dim_F {}$\leanname{SwapRigidity.swapCode}$(A, B) = d^2$.
\end{enumerate}
\end{proposition}
\begin{outline}
In~(a), the intervals used in $A$ create at most $2\sum_j M_A(j)$ endpoint cuts around $\Z_d$,
partitioning $\mathrm{Fin}\,d$ into at most $\beta_A$ contiguous classes on each of which every $A_a$
acts as a fixed linear combination of rotations. Hence $\sum_a A_a X B_a^{\T} = X^{\T}$ defines a
$\beta_A$-class-uniform code, giving $2\dim_F \le d^2 + d\beta_A$ by
\cite[Theorem~6.4]{SeventeenthPaper2026}. Passing to $Y = X^{\T}$ exchanges $A$ and $B$ via
$Y^{\T} = \sum_a B_a Y A_a^{\T}$, giving $2\dim_F \le d^2 + d\beta_B$ as well. Part~(b) divides by
$2d^2$. In~(c), $\sum_{j=0}^{m-1} 2^{j+1} = 2d - 2 \ge d$ for $m \ge 1$, and aligned dyadic intervals
down to length $2^0 = 1$ span all diagonal projectors
$E_{r,r}$ and hence all matrix units $E_{r,c} = E_{r,r}\sigma^{c-r}$ (where $(\sigma v)_a = v_{a+1}$).
\end{outline}

Note the scope of \Cref{prop:multiscale-interval}: the bound uses only the sparser of the two sides
rather than any two-sided interaction between $A$ and $B$, carries no matching sparse lower bound,
and becomes vacuous ($\ge d^2$) as soon as $\min(\beta_A, \beta_B) \ge d$---in particular at the full
profile $M_{\mathrm{full}}(j) = 2^{j+1}$ for $m \ge 2$, where $\beta = 4d - 4 > d$.

To exploit sparsity on \emph{both} sides simultaneously in a two-scale (coarse--fine) setting over an
arbitrary field $F$, suppose the intervals used on side $A$ separate into a ``fine'' coordinate set
$S_A \subseteq \mathrm{Fin}\,d$ of cardinality $s_A = \abs{S_A}$ and $c_A$ pairwise-disjoint nonempty
coarse cyclic intervals $\mathrm{coarse}_A : \mathrm{Fin}\,c_A \to \mathrm{Finset}(\mathrm{Fin}\,d)$
partitioning the complement $L_A := \mathrm{Fin}\,d \setminus S_A$, so that every $A_a$ lies in the
$F$-span of $\{\lean{intervalMaskRotMatrix}(F, I, e)\}$ with $I \subseteq S_A$,
$I = \mathrm{coarse}_A(i)$, or $I = \mathrm{Fin}\,d$ (\lean{HasCoarseFineIntervalDecomposition}), and
similarly $(S_B, c_B)$ on side $B$ ($s_B = \abs{S_B}$, $L_B = \mathrm{Fin}\,d \setminus S_B$). Note
that \lean{HasCoarseFineIntervalDecomposition} allows only masks supported in $S$, single coarse
intervals partitioning $\mathrm{Fin}\,d \setminus S$, or the full set $\mathrm{Fin}\,d$: nested
dyadic families with three or more active scales force intermediate-scale intervals into $S$, so
\Cref{thm:two-sided-coarse-fine} below is a bound for two-scale (coarse--fine) families rather than a
general multi-scale profile bound for \cite[Problem~9.3]{SeventeenthPaper2026}. For any
$S \subseteq \mathrm{Fin}\,d$, write
$\dot{m}(S) := \lean{cyclicRunEndpointCount}(S) = \abs{\{r \in S : (r + 1) \bmod d \notin S\}}$
and $m(S) := \lean{cyclicRunCount}(S)$ ($= 1$ if $S = \mathrm{Fin}\,d$ and $\dot{m}(S)$ otherwise, so
$\dot{m}(S) \le m(S)$), and set (\lean{theorem3UnsymmetrizedBound})
\begin{align*}
  B_{\mathrm{cf}}(d;\, s_A, m_A, c_A;\, s_B, m_B, c_B)
  &\;:=\;
  \frac{1}{2}\,d^2
  + \frac{1}{4}\,s_B(d - s_A)
  + \frac{1}{2}\,s_A s_B \\
  &\quad + \frac{1}{2}\,c_B d
  + \frac{1}{4}\,c_A s_B
  + \frac{1}{2}\,s_B m_A
  + \frac{1}{4}\,m_B(d - s_A).
\end{align*}

\begin{theorem}[Characteristic-free two-sided coarse--fine bound for interval-rotation swap codes]\label{thm:two-sided-coarse-fine}
\formal{ExtensionTwoSidedProfileProofs.lean}{extension_Theorem3_TwoSidedCoarseFine,swapCode_finrank_le_twoSidedCoarseFine}
Let $F$ be \emph{any} field (of arbitrary characteristic), let $d \ge 1$, let $\alpha$ be a finite
index type, and let $A, B : \alpha \to \Mat_{d\times d}(F)$ satisfy
\lean{HasCoarseFineIntervalDecomposition} $(S_A, c_A, A)$ and
\lean{HasCoarseFineIntervalDecomposition} $(S_B, c_B, B)$ with $s_A = \abs{S_A}$ and
$s_B = \abs{S_B}$. Then with either $(m_A, m_B) = \bigl(\dot{m}(S_A), \dot{m}(S_B)\bigr)$ or
$(m_A, m_B) = \bigl(m(S_A), m(S_B)\bigr)$,
\begin{align*}
  \dim_F \leanname{SwapRigidity.swapCode}(A, B)
  &\;\le\;
  \min\!\Bigl(
    B_{\mathrm{cf}}(d;\, s_A, m_A, c_A;\, s_B, m_B, c_B), \\
  &\qquad\qquad
    B_{\mathrm{cf}}(d;\, s_B, m_B, c_B;\, s_A, m_A, c_A)
  \Bigr).
\end{align*}
\end{theorem}
\begin{outline}
Identify $\mathrm{Fin}\,d$ with $\Z_d$ and let $K(X) := \sum_{a \in \alpha} A_a X B_a^{\T}$, so
$\mathcal{X} := \leanname{SwapRigidity.swapCode}(A, B) = \{X : X^{\T} = K(X)\}$.
\begin{enumerate}[label=\arabic*.,leftmargin=1.5em]
\item \emph{Corner perturbation}: subtracting the fine-by-fine corner $X \mapsto P_{S_A} K(X) P_{S_B}$
  (where $(P_S v)_a = \mathbf{1}_{a \in S}\,v_a$, image dimension $\le s_A s_B$) gives
  $\dim_F \mathcal{X} \le s_A s_B + \dim_F \mathcal{Y}$, where
  $\mathcal{Y} := \{X : X^{\T} = K(X) - P_{S_A} K(X) P_{S_B}\}$. Every $X \in \mathcal{Y}$ has
  $(X^{\T})_{r,c} = 0$ on $S_A \times S_B$, i.e.\ $X_{b,a} = 0$ for all $(b,a) \in S_B \times S_A$.
\item \emph{Column-shift Grassmann step}: for $\sigma_{\mathrm{col}}(X)_{b,a} := X_{b, a+1}$, the sum
  $\mathcal{Y} + \sigma_{\mathrm{col}}^{-1}(\mathcal{Y})$ vanishes on
  $S_B \times \{a \in S_A : a - 1 \in S_A\}$ of cardinality $s_B(s_A - \dot{m}(S_A))$, so
  $2\dim_F \mathcal{Y} + s_B(s_A - \dot{m}(S_A)) \le d^2 + \dim_F \mathcal{Y}_1$ where
  $\mathcal{Y}_1 := \mathcal{Y} \cap \sigma_{\mathrm{col}}^{-1}(\mathcal{Y})$.
\item \emph{$B$-side coarse anchors}: along each $B$-coarse interval $\mathrm{coarse}_B(j)$ with start
  $\mathrm{start}_B(j)$, shift covariance of the $B_a$ forces $X_{b+1, a} = X_{b, a+1}$ across all
  $a \in \Z_d$ for every $X \in \mathcal{Y}_1$ and every interior step of $\mathrm{coarse}_B(j)$.
  Setting the $c_B d$ anchor entries $X_{\mathrm{start}_B(j), a} = 0$ defines $\mathcal{Z}_0 \le
  \mathcal{Y}_1$ with $\dim_F \mathcal{Y}_1 \le c_B d + \dim_F \mathcal{Z}_0$, and every
  $X \in \mathcal{Z}_0$ vanishes outside $S_B \times L_A$.
\item \emph{Row-shift Grassmann step}: for $\sigma_{\mathrm{row}}(X)_{b,a} := X_{b+1, a}$, the sum
  $\mathcal{Z}_0 + \sigma_{\mathrm{row}}^{-1}(\mathcal{Z}_0)$ is supported inside
  $(S_B \cup (S_B + 1)) \times L_A$ of size at most $(s_B + \dot{m}(S_B))(d - s_A)$, while on
  $\mathcal{Z}_1 := \mathcal{Z}_0 \cap \sigma_{\mathrm{row}}^{-1}(\mathcal{Z}_0)$ shift covariance of
  the $A_a$ along each $A$-coarse interval forces $X_{b, \mathrm{start}_A(i) + u} =
  X_{b + u, \mathrm{start}_A(i)}$, injecting $\mathcal{Z}_1$ into $F^{\mathrm{Fin}\,c_A \times S_B}$ of
  dimension $c_A s_B$. Hence $2\dim_F \mathcal{Z}_0 \le (s_B + \dot{m}(S_B))(d - s_A) + c_A s_B$.
\end{enumerate}
Combining steps~1--4 gives $\dim_F \mathcal{X} \le B_{\mathrm{cf}}(d;\, s_A, \dot{m}(S_A), c_A;\,
s_B, \dot{m}(S_B), c_B)$; transposing $Y = X^{\T}$ swaps $(A, S_A, c_A)$ with $(B, S_B, c_B)$, and
monotonicity in $(m_A, m_B)$ gives the $m(S)$ form since $d - s_A \ge 0$.
\end{outline}

On the concrete family \lean{InContiguousCoarseFineSpan} generated on side $A$ by $s_A$ contiguous
singleton masks $\{0\}, \dots, \{s_A - 1\}$ (\lean{contiguousFineSet}, $m_A = 1$) together with one
complementary coarse interval $\{s_A, \dots, d - 1\}$ (\lean{contiguousCoarseComplement}, $c_A = 1$)
and full rotations---and analogously $s_B$ contiguous singletons plus one coarse complement on side
$B$---applying Theorem~TP's interval clause (\cite[Theorem~6.4]{SeventeenthPaper2026}, $(d^2 + Md)/2$)
to one side at a time partitions $\mathrm{Fin}\,d$ into $M = s_A + 1$ cyclic intervals on side $A$ (so
$\sum_a A_a X B_a^{\T} = X^{\T}$ makes $X^{\T}$ an $(s_A + 1)$-class-uniform code) and $s_B + 1$
cyclic intervals on side $B$ (passing to $Y = X^{\T}$). That one-sided Theorem~TP baseline gives
$(d^2 + (\min(s_A, s_B) + 1)d)/2$, with leading coefficient $\frac{1}{2} + \frac{1}{2}\min(\alpha,
\beta)$ ($\alpha = s_A / d$, $\beta = s_B / d$), whereas \Cref{thm:two-sided-coarse-fine} achieves a
strictly smaller leading coefficient for all $\alpha, \beta \in (0, 1)$ (note that at small $d$ or
$s = 1$, Theorem~TP can be tighter: on $s_A = s_B = s$,
$\frac{1}{2}(d^2 + (s+1)d) - B_{\mathrm{cf}}(d; s,1,1; s,1,1) = \frac{1}{4}\bigl(s(d - s) - 2s - d\bigr)$,
and if $S$ is scattered with $\dot{m}(S) = \Theta(d)$, such as the even residues, $B_{\mathrm{cf}}$
can exceed $d^2$):

\begin{corollary}[Leading-coefficient win over one-sided Theorem~TP and $\frac{11}{16}d^2 + d$ benchmark]\label{cor:two-sided-win}
\partformal{ExtensionTwoSidedProfileProofs.lean}{extension_Theorem3_ConcreteWinOverTP,swapCode_finrank_le_contiguousCoarseFine}{applying \cite[Theorem~6.4]{SeventeenthPaper2026} to \lean{InContiguousCoarseFineSpan} to obtain the one-sided baseline $(d^2 + (\min(s_A, s_B) + 1)d)/2$ (part~(a) in Lean is the real-coefficient inequality).}
\begin{enumerate}[label=(\alph*)]
\item For all real $\alpha, \beta \in (0, 1)$,
  \[
    \frac{1}{2} \;+\; \frac{1}{4}\min\bigl(\beta(1 + \alpha),\; \alpha(1 + \beta)\bigr)
    \;<\;
    \frac{1}{2} \;+\; \frac{1}{2}\min(\alpha, \beta).
  \]
\item For every field $F$, every finite index type $\alpha_0$, every $d \ge 2$, every
  $1 \le s_A, s_B < d$ ($\alpha := s_A / d$, $\beta := s_B / d$), and all
  $A, B : \alpha_0 \to \Mat_{d\times d}(F)$ satisfying
  \lean{InContiguousCoarseFineSpan} $(s_A, A)$ and \lean{InContiguousCoarseFineSpan} $(s_B, B)$, both
  sides satisfy \lean{HasCoarseFineIntervalDecomposition} with $c_A = c_B = m_A = m_B = 1$, and
  \[
    \dim_F \leanname{SwapRigidity.swapCode}(A, B)
    \;\le\;
    \left(\frac{1}{2} + \frac{1}{4}\min\bigl(\beta(1 + \alpha),\; \alpha(1 + \beta)\bigr)\right) d^2
    \;+\; \frac{5}{4}\,d.
  \]
\item At the half-fine benchmark $d = 2k$ ($k \ge 1$) with $s_A = s_B = k$ ($\alpha = \beta = 1/2$),
  one has $B_{\mathrm{cf}}(2k;\, k, 1, 1;\, k, 1, 1) = \frac{11}{16}\,d^2 + d$
  ($\frac{11}{16} = 0.6875 < \frac{3}{4}$), so every pair $A, B$ in
  \lean{InContiguousCoarseFineSpan} $(k)$ satisfies
  \[
    \dim_F \leanname{SwapRigidity.swapCode}(A, B) \;\le\; \frac{11}{16}\,d^2 + d,
  \]
  which is strictly less than $\frac{3}{4}\,d^2$ for all even $d = 2k \ge 18$ ($k \ge 9$).
\end{enumerate}
\end{corollary}
\begin{outline}[Outline of the formal proof and proof of the one-sided baseline]
For the one-sided baseline (proved in text), on any $A$ in \lean{InContiguousCoarseFineSpan}$(s_A)$,
every $A_a$ acts on each of the $s_A + 1$ cyclic intervals $\{0\}, \dots, \{s_A - 1\},
\{s_A, \dots, d - 1\}$ as a fixed linear combination of rotations, so $\sum_a A_a X B_a^{\T} = X^{\T}$
places $X^{\T}$ in an $(s_A + 1)$-interval class-uniform code of dimension at most
$(d^2 + (s_A + 1)d)/2$ by \cite[Theorem~6.4]{SeventeenthPaper2026} (and $(d^2 + (s_B + 1)d)/2$ on
$X$ by symmetry). In~(a), assuming $\alpha \le \beta$ by symmetry gives
$\min\bigl(\beta(1+\alpha), \alpha(1+\beta)\bigr) \le \alpha(1+\beta) < 2\alpha = 2\min(\alpha,\beta)$
since $\alpha > 0$ and $1 + \beta < 2$. In~(b), $\lean{contiguousFineSet}(s)$ has cardinality $s$ and
a single cyclic run ($m = 1$) while $\lean{contiguousCoarseComplement}(s)$ is a single cyclic interval
of length $d - s$ ($c = 1$), and substituting $c_A = c_B = m_A = m_B = 1$ into
\Cref{thm:two-sided-coarse-fine} yields
$\frac{1}{2}d^2 + \frac{1}{4}s_B d + \frac{1}{4}s_A s_B + \frac{1}{2}d + \frac{3}{4}s_B +
\frac{1}{4}(d - s_A) \le (\frac{1}{2} + \frac{1}{4}\beta(1+\alpha))d^2 + \frac{5}{4}d$ when
$s_B \le s_A$ (and symmetrically when $s_A \le s_B$). Part~(c) evaluates
$B_{\mathrm{cf}}(2k;\, k,1,1;\, k,1,1) = 2k^2 + \frac{1}{4}k^2 + \frac{1}{2}k^2 + k + \frac{1}{4}k +
\frac{1}{2}k + \frac{1}{4}k = \frac{11}{4}k^2 + 2k = \frac{11}{16}d^2 + d$, which satisfies
$\frac{11}{16}d^2 + d < \frac{3}{4}d^2 \iff d > 16 \iff k \ge 9$.
\end{outline}

\begin{remark}[Numerical: rank-minimization search on small two-sided interval partitions]\label{rem:two-sided-numerical}
A generic random operator $K \in \mathcal{U}(\mathcal{I}) \otimes \mathcal{U}(\mathcal{J})$ in the
two-sided interval-rotation span has $\dim\ker(\mathrm{sw} - K) = 0$. To probe maximum swap-code
dimensions over $\R$ at $d \in \{5, 6, 7, 8\}$ (reproducible via
\leanfile{tools/verify_two_sided_svd.py}; not formalized, and no proof in this paper uses this
computation), we ran Nesterov-accelerated alternating projections with adaptive restart between
$\mathcal{U}(\mathcal{I}) \otimes \mathcal{U}(\mathcal{J})$ and bounded-rank manifolds
$\{K : \rank(\mathrm{sw} - K) \le d^2 - D\}$ via truncated SVD across random restarts. Writing
$D_{\diamond}(\mathcal{I}, \mathcal{J}) := \bigl\lfloor \frac{1}{2}\bigl(d^2 + \sum_{i, j}
\min(\abs{I_i}, \abs{J_j})\bigr)\bigr\rfloor$, the search found no witness at
$D = D_{\diamond} + 1$ on any tested partition. Bounded-norm operators converged to machine precision
($\le 10^{-11}$) at $D = D_{\diamond}$ for the contiguous fine-singleton + coarse-complement
partitions $(d, s) \in \{(5, 2), (6, 3), (7, 4), (8, 4)\}$ (including the half-density cases
$s = d/2$, where $D_{\diamond} = \frac{5}{8}d^2 + O(d)$), for $d = 6$ with $s_A = 3$ and $s_B = 4$,
and for the $d = 6$ two-run partition $\{0\}, \{1,2\}, \{3\}, \{4,5\}$; by contrast, $D_{\diamond}$
was not reached by the budgeted search for $(d, s) = (7, 3)$ ($D_{\diamond} = 34$), $(d, s) = (8, 3)$
($D_{\diamond} = 42$), and the $d = 7$ partition $1 + 1 + 5$ against $2 + 2 + 3$
($D_{\diamond} = 31$), where the witnesses found have dimensions $32$, $40$, and $27$. Whether
$D_{\diamond}(\mathcal{I}, \mathcal{J})$ is an upper bound for all two-sided interval partitions
remains open.
\end{remark}

By contrast, if ``scale-$j$ dyadic mask'' means the periodic Walsh/Rademacher bit projector
$P_{W_j, \mathrm{bit}}$ onto $\{r \in \mathrm{Fin}(2^m) : \mathrm{bit}_j(r) = \mathrm{bit}\}$
($\mathrm{bit} \in \{0, 1\}$; \lean{walshBitSet},
\lean{walshDigitMaskRotMatrix}, \lean{IsPeriodicScaleDyadicMask}), then a single mask per scale has
exponentially many boundaries at fine scales, so one-sided endpoint refinement becomes vacuous even
though only two periodic masks per scale are used:

\begin{proposition}[Cyclic boundaries and operator-span dimension for periodic Walsh masks]\label{prop:walsh-dichotomy}
\formal{SwapCodeAnswersProofs.lean}{walshDigitMask_in_periodicDyadicSpan}
\formal{DualityAndSwapConsequences.lean}{walshBitSet_cyclicBoundaryCount,isCyclicInterval_cyclicBoundaryCount_le_two,isPeriodicScaleDyadicMask_cyclicBoundaryCount,walshBitSet_joint_separates_points,walshDigitMaskRotSpan_finrank_le,bonus2_WalshDichotomyStructuralProperties}
Let $m \ge 0$ and $d = 2^m$.
\begin{enumerate}[label=(\alph*)]
\item Every cyclic interval $I \subseteq \mathrm{Fin}\,d$ ($d \ge 1$) satisfies
  $\lean{cyclicBoundaryCount}(I) \le 2$, whereas for each $j \in \mathrm{Fin}\,m$, both
  $W_j = \lean{walshBitSet}(m, j)$ and every periodic scale-$j$ mask
  (\lean{IsPeriodicScaleDyadicMask}) have $\lean{cyclicBoundaryCount} = 2^{j+1}$ (at least $4$ once
  $j \ge 1$).
\item Across $j \in \mathrm{Fin}\,m$, the $m$ Walsh bit sets $W_0, \dots, W_{m-1}$ jointly separate
  $\mathrm{Fin}(2^m)$ into $2^m$ singleton atoms.
\item Over every field $F$, every family in the span of
  $\{\lean{walshDigitMaskRotMatrix}(F, j, \mathrm{bit}, s) : j \in \mathrm{Fin}\,m,\;
  \mathrm{bit} \in \{0,1\},\; s \in \mathrm{Fin}(2^m)\}$ belongs to
  \lean{InMultiscalePeriodicDyadicRotSpan} with constant profile $M(j) = 2$ per scale
  ($\sum_{j=0}^{m-1} M(j) = 2m = 2\log_2 d$), and the entire Walsh-masked rotation span has dimension
  at most $(m + 1)2^m$ over $F$.
\end{enumerate}
\end{proposition}
\begin{outline}
In~(a), $W_j$ alternates blocks of $2^{m-1-j}$ zeros and $2^{m-1-j}$ ones $2^j$ times around
$\mathrm{Fin}(2^m)$, producing $2^{j+1}$ boundary transitions, and any periodic scale-$j$ mask is a
cyclic translate of $W_j$ or its complement. Part~(b) recovers $r \in \mathrm{Fin}(2^m)$ from its
$m$ binary digits. In~(c), the two bit values $\mathrm{bit} \in \{0,1\}$ at scale $j$ give two
periodic scale-$j$ masks ($M(j) = 2$), while $P_{W_j, 0}\sigma^s + P_{W_j, 1}\sigma^s = \sigma^s$
shows that the span is generated by the $2^m$ pure rotations together with the $m \cdot 2^m$
operators $P_{W_j, 1}\sigma^s$.
\end{outline}

Whether $\dim_F \leanname{SwapRigidity.swapCode}(A, B)$ is bounded away from $d^2$ when both $A$ and
$B$ lie in the at most $(m+1)2^m$-dimensional Walsh-masked rotation span (a periodic variant of
\cite[Problem~9.3]{SeventeenthPaper2026}) remains open.
